\chapter[Introdução]{Introdução}

\section{Contextualização}

O Trabalho de Conclusão de Curso (TCC) consiste em um trabalho acadêmico desenvolvido ao final de uma formação. De acordo com a Associação Brasileira de Normas Técnicas (ABNT), por meio da NBR 14724, o TCC é um documento que representa o resultado de um estudo, no qual o estudante deve expressar conhecimento do assunto escolhido, obrigatoriamente vinculado à disciplina, módulo, curso ou programa, e elaborado sob a coordenação de
um orientador \cite{ABNT2024}.

% O desenvolvimento do TCC envolve diferentes etapas e a participação de estudantes, orientadores e da instituição de ensino. O acompanhamento do cronograma e a interação entre orientador e orientando constituem aspectos desse processo, conforme discutido por \citeonline{Viana2010}.

% A relação entre orientador e orientando pode apresentar
% dificuldades ao longo do desenvolvimento do trabalho.
% \citeonline{Viana2010} destacam que o cumprimento do
% cronograma é uma responsabilidade compartilhada entre
% ambos, envolvendo a realização das atividades previstas
% e o acompanhamento das orientações. Entretanto, as autoras
% identificam dificuldades nesse processo, como a necessidade
% de alguns orientadores buscarem contato com seus orientandos e a sobrecarga de trabalho de determinados docentes, que pode comprometer a disponibilidade para o atendimento. De maneira semelhante, \citeonline{LeiteFilho2006} apontam a falta de contato entre orientador e orientando como uma das principais dificuldades dessa relação, frequentemente associada à falta de tempo dos envolvidos.

% Na Faculdade de Ciências e Tecnologias em Engenharia (FCTE) da Universidade de Brasília (UnB), o TCC é regulamentado pela Resolução do Conselho da FCTE nº 0002/2022, que estabelece sua organização em dois componentes curriculares. O TCC 1 possui caráter qualificatório, com o objetivo de avaliar a aptidão e a maturidade do discente e do tema em sua área de investigação. Já o TCC 2 corresponde à versão definitiva do trabalho, destinada à sua apresentação e validação. Ambos os componentes são avaliados por bancas examinadoras específicas.
% Após a avaliação, o resultado é registrado em ata e, no caso do
% TCC 2 aprovado, a versão definitiva é encaminhada para
% disponibilização na Biblioteca Digital de Monografias (BDM)
% \cite{UnB2022}.

A orientação acadêmica constitui uma parte importante do desenvolvimento do trabalho e pode apresentar dificuldades ao longo desse processo.\citeonline{Viana2010} descrevem a orientação acadêmica como uma relação pedagógica estabelecida entre orientador e orientando e apontam a existência de contribuições e fragilidades relacionadas a ambos os envolvidos. As autoras destacam que, por se tratar de uma relação dinâmica entre indivíduos com diferentes características pessoais e profissionais, essa interação pode se estabelecer de maneiras distintas, o que reforça a importância do diálogo entre as partes na construção do trabalho acadêmico.

Além do diálogo estabelecido durante a orientação, o acompanhamento do cronograma e a participação nas etapas também fazem parte desse processo. \citeonline{Viana2010} destacam que o cumprimento do cronograma é uma responsabilidade compartilhada entre orientador e orientando. As autoras identificam, entretanto, dificuldades nesse acompanhamento, tanto pela necessidade de alguns orientadores buscarem ativamente seus orientandos quanto pela sobrecarga de trabalho de alguns docentes, que pode comprometer o atendimento aos alunos. Semelhantemente, \citeonline{LeiteFilho2006} apontam a ausência de contato, decorrente da falta de tempo de ambas as partes, como o principal fator que dificulta a relação, e observam que a frequência dos encontros costuma depender da iniciativa do orientando.

% Além do orientador e do orientando, a coordenação também faz parte do processo, aprovando a composição das bancas, acompanhando o andamento dos trabalhos, controlando os prazos e reunindo as atas e menções ao final de cada semestre. Embora o regulamento da FCTE faça referência a um sistema eletrônico próprio para a matrícula nos componentes, a divulgação de menções, o registro das atas e a emissão de comprovante de participação ou orientação de banca, esses sistemas se limitam a etapas administrativas pontuais e não acompanham o trabalho enquanto ele está sendo desenvolvido. Nesse sentido, reuniões, comentários sobre versões do texto e ajustes de cronograma costumam ser registrados em diferentes meios de comunicação, como email e aplicativos de mensagem, sem um histórico centralizado que permita aos envolvidos acompanhar o progresso do trabalho.

% Dessa forma, o desenvolvimento do TCC envolve diferentes atividades, etapas e participantes, não se restringindo apenas à elaboração do trabalho acadêmico. O processo demanda acompanhamento ao longo de sua realização e envolve a interação entre estudantes, orientadores e a própria instituição de ensino. Nesse contexto,  Assim, este trabalho tem como foco a organização e o acompanhamento das diferentes etapas que compõem o ciclo de vida do TCC.

Dessa forma, o desenvolvimento do TCC não se restringe à elaboração do documento acadêmico. O processo demanda acompanhamento ao longo de sua realização e envolve a interação entre estudantes, orientadores e a própria instituição de ensino. Nesse contexto, este trabalho tem como foco a gestão e o acompanhamento do ciclo de vida do Trabalho de Conclusão de Curso.


\section{Justificativa}

O uso de tecnologias digitais tem se tornado cada vez mais presente no ambiente acadêmico, especialmente a partir das transformações ocorridas durante a pandemia de COVID-19. Mesmo após o retorno das atividades presenciais, práticas digitais incorporadas nesse período permaneceram presentes no ensino superior, sendo utilizadas em diferentes atividades acadêmicas e profissionais \cite{Carneiro2026}.

Nesse contexto, a utilização de ambientes virtuais passou a envolver não apenas o acesso a conteúdos, mas também atividades relacionadas à organização, à gestão de tarefas e ao acompanhamento das interações acadêmicas  \cite{Carneiro2026}. Essa realidade demonstra a aplicação dos recursos digitais no suporte a diferentes processos acadêmicos, criando possibilidades para sua utilização também na organização e no acompanhamento de atividades que envolvem diferentes etapas e participantes.

Nesse panorama, o processo de elaboração do TCC apresenta características que favorecem a utilização desse suporte. Visto que envolve etapas sequenciais, prazos definidos, entregas intermediárias e a participação de diferentes atores, seu acompanhamento depende da organização e atualização de informações que, atualmente, permanecem dispersas. Nesse sentido, a centralização dessas informações em um único ambiente pode contribuir para uma visão mais organizada do processo, facilitando o acompanhamento das etapas e das responsabilidades de cada um dos envolvidos.

% Apesar da existência de ferramentas voltadas para etapas específicas da gestão do TCC, o processo de desenvolvimento envolve informações e tarefas que vão além desses registros pontuais. A escolha do orientador, o planejamento das atividades, o acompanhamento de prazos, o registro de reuniões, o compartilhamento de documentos e o controle das entregas frequentemente ocorrem de maneira descentralizada. 

Diante desse cenário, faz-se necessária uma solução digital voltada especificamente para o acompanhamento do ciclo de vida do TCC. Um ambiente centralizado pode reunir informações relacionadas ao trabalho, como o andamento das etapas, os registros das reuniões e os documentos produzidos ao longo do processo, facilitando a comunicação entre o estudante e o orientador e oferecendo à coordenação uma visão ampla dos trabalhos em andamento e de suas respectivas situações.

% Dado o exposto, a escolha deste tema justifica-se pela necessidade de tornar o processo de desenvolvimento e acompanhamento dos trabalhos de Conclusão de Curso mais organizado, integrado e acessível aos envolvidos. A existência de informações, atividades e documentos distribuídos em diferentes meios pode dificultar tanto o acompanhamento realizado por alunos e orientadores quanto a visão da coordenação sobre o andamento dos trabalhos. Portanto, o desenvolvimento de um sistema voltado especificamente para esse contexto apresenta-se como uma oportunidade de reunir, em um único ambiente, recursos que auxiliem na organização das etapas, no acompanhamento de prazos e atividades, no registro das interações e na gestão das informações produzidas ao longo do TCC. Assim, a escolha pelo desenvolvimento dessa solução busca contribuir para uma gestão melhor estruturada do processo, atendendo às necessidades identificadas ao longo do ciclo de vida do Trabalho de Conclusão de Curso.


\section{Questão da Pesquisa}

Diante das dificuldades encontradas, este trabalho busca responder à seguinte questão:

\textbf{Como um sistema integrado pode contribuir para a gestão e o acompanhamento do ciclo de vida do Trabalho de Conclusão de Curso, considerando as necessidades de alunos, orientadores e coordenação?}

Nesse contexto, propõe-se o desenvolvimento de uma aplicação web que facilite a gestão e o acompanhamento das etapas do TCC, considerando as necessidades dos envolvidos.

\section{Objetivos}

Esta seção apresenta o objetivo geral e os objetivos específicos deste trabalho.

\subsection{Objetivo Geral}

Desenvolver uma aplicação web que centralize e facilite a gestão e o acompanhamento do ciclo de vida do TCC, atendendo às necessidades de alunos, orientadores e coordenação.

\subsection{Objetivos Específicos}

Os objetivos específicos descrevem as etapas necessárias para o desenvolvimento da solução proposta, considerando as necessidades dos usuários, os requisitos do sistema e as práticas de Engenharia de Software adotadas ao longo do projeto. Nesse sentido, foram definidos os seguintes objetivos específicos:

\begin{enumerate}
    \item Levantar as necessidades e dificuldades de estudantes, orientadores e coordenação no processo de TCC por meio da aplicação de questionários.
    \item Identificar e definir os requisitos do sistema a partir das necessidades levantadas.
    \item Projetar a interface e os fluxos de interação do sistema.
    \item Desenvolver a interface web responsiva do sistema.
    \item Desenvolver a API responsável pelas regras de negócio e pela comunicação entre os componentes da aplicação.
    \item Implementar o armazenamento e o gerenciamento dos dados da aplicação.
    \item Avaliar a solução desenvolvida com os usuários envolvidos no processo de TCC.
\end{enumerate}

Os objetivos específicos foram divididos em duas etapas, sendo os quatro primeiros realizados na primeira entrega deste trabalho, enquanto os objetivos de 5 a 7 serão desenvolvidos na segunda etapa.


\section{Organização da Monografia}



\section{Uso de Inteligência Artificial}

Durante a realização deste trabalho, ferramentas de Inteligência Artificial (IA) serão utilizadas como apoio à elaboração do documento acadêmico e ao desenvolvimento da aplicação. Na produção textual, seu uso será direcionado à revisão ortográfica e à melhoria da clareza e organização do texto. 

No desenvolvimento da aplicação, a IA será utilizada na construção da interface web, incluindo a geração e o aprimoramento de componentes e funcionalidades do frontend. O código produzido será revisado e validado pelos autores, que permanecerão responsáveis pelas decisões técnicas e pela implementação da solução.